\chapter*{Project Objective}
\thispagestyle{fancy}
\addcontentsline{toc}{chapter}{Project Objective}
\begingroup
\setlength{\emergencystretch}{2em}

The objective of this project is to establish a traceable antenna-engineering workflow
for an inset-fed rectangular microstrip patch operating near 915~MHz. The work connects
first-order transmission-line synthesis to a parameterized CST model and evaluates the
final input match, near-field behavior, and broadside radiation characteristics.

The project has four principal technical objectives:

\begin{enumerate}
    \item Calculate the initial patch width and length from the design frequency,
    substrate permittivity, substrate thickness, effective dielectric constant, and
    fringing-field correction.

    \item Define a practical inset-fed geometry using a documented 50-$\Omega$
    microstrip width and a first-order inset-depth estimate based on the patch input
    resistance.

    \item Construct and tune the antenna in CST Studio Suite, then quantify the final
    resonant frequency and reflection coefficient relative to the configured port.

    \item Evaluate the preserved electric-field, magnetic-field, 3D directivity, and
    principal-plane radiation-pattern evidence at the 915~MHz design frequency.
\end{enumerate}

A secondary objective is to distinguish exact displayed markers, values read directly
from CST information panels, analytical calculations, and qualitative field evidence.
The report also identifies the numerical and experimental work still required before
hardware-performance claims would be justified.

The final outcome is a concise, university-style public engineering report consistent
with the other antenna projects in this portfolio. All reported performance is
analytical or simulated; no fabricated or measured antenna result is claimed.
\endgroup
